\section{Нефункціональні вимоги та системні обмеження}\label{sec:nefunktsional}

Усі нефункціональні вимоги мають ідентифікатори формату \texttt{REQ-NF-xxx} і зосереджені на якісних характеристиках надійності, безпеки та транзакційної цілісності даних.

\subsection{Транзакційна атомарність та надійність даних}\label{sec:nefunktsional-atomarnist}

\begin{itemize}[leftmargin=*,labelindent=0pt]
    \item \textbf{REQ-NF-001 (Транзакційна атомарність створення бронювання):} Система повинна гарантувати сувору атомарність (властивість ACID) для операції створення бронювання. Будь-який збій на проміжному етапі повинен повністю відкочувати транзакцію до попереднього узгодженого стану (принцип «все або нічого»)~\sked{K5,K30}.
    \textit{Умова верифікації:} Штучне переривання транзакції запису в базу даних не залишає частково заброньованих місць або бронювань без прив'язки до місць.

    \item \textbf{REQ-NF-002 (Персистентність та збереження цілісності):} Серверна частина повинна гарантувати персистентне збереження та цілісність усіх сутностей (користувачі, фільми, зали, сеанси, бронювання). Рестарт системи чи падіння окремих процесів не повинні призводити до втрати узгодженості або дублювання даних~\sked{K40,K41}.
    \textit{Умова верифікації:} Автоматичний тест перезавантаження сервера підтверджує 100\% збереження всіх зафіксованих перед зупинкою бронювань.
\end{itemize}

\subsection{Безпека та захист від несанкціонованого доступу}\label{sec:nefunktsional-bezpeka}

\begin{itemize}[leftmargin=*,labelindent=0pt]
    \item \textbf{REQ-NF-003 (Серверний контроль прав RBAC):} Усі операції зміни даних та доступу до приватних ресурсів повинні перевірятися на рівні бізнес-логіки сервера. Клієнтські перевірки в браузері розглядаються виключно як засіб підвищення зручності інтерфейсу, але не як засіб безпеки~\sked{K40}.
    \textit{Умова верифікації:} Запити в обхід вебклієнта (через прямі виклики API) суворо підпорядковуються серверним правилам авторизації.

    \item \textbf{REQ-NF-004 (Ізоляція даних між користувачами):} Система повинна унеможливлювати витік інформації про бронювання іншого клієнта чи маніпулювання чужими бронюваннями (захист від атак типу Insecure Direct Object References, IDOR)~\sked{K24}.
    \textit{Умова верифікації:} Зміна ID бронювання в тілі або параметрах запиту на ID чужого бронювання повертає помилку доступу.

    \item \textbf{REQ-NF-005 (Неможливість самопризначення ролі адміністратора):} Процедура реєстрації звичайного користувача повинна створювати обліковий запис виключно з роллю «Користувач». Спроби передачі адміністративних прапорців у запиті реєстрації повинні ігноруватися або відхилятися~\sked{K27}.
    \textit{Умова верифікації:} Перевірка створеного запису в базі даних після реєстрації із запитом, що містить поле \texttt{role: "admin"}, підтверджує присвоєння ролі звичайного користувача.
\end{itemize}

\subsection{Обробка конкурентних запитів та продуктивність}\label{sec:nefunktsional-konkurentnist}

\begin{itemize}[leftmargin=*,labelindent=0pt]
    \item \textbf{REQ-NF-006 (Коректність при конкурентному навантаженні):} Відповідно до узгоджених вимог, для першої версії не встановлюються штучні числові пороги часу відповіді чи кількості одночасних користувачів. Головною вимогою до продуктивності є якісна здатність системи надійно опрацьовувати одночасні запити до одного й того самого сеансу без взаємних блокувань (дедлоків), деградації цілісності бази даних та гарантувати, що одне місце ні за яких умов не буде заброньоване двома різними клієнтами~\sked{K29,K43}.
    \textit{Умова верифікації:} Стрес-тест конкурентного доступу (100 паралельних запитів на одні й ті самі місця) підтверджує відсутність подвійних бронювань та завершення транзакцій без дедлоків.
\end{itemize}
